% Chapter 8 — Ingestion, Corporate Actions, and the Market Data Operator (NAZAROV, 8.5 pp).
% Discharges Constitution v1.2 §9. Carries G1 (four mechanisms + spin-off/elective-merger
% half-page examples), G3 (W1-W4), G5 (cum/ex). Declares the boundary-convention slot Ch.12 consumes.
\section{Ingestion, Corporate Actions, and the Market Data Operator}\label{ch:marketdata}

This chapter discharges constitution clauses C-9.1--C-9.3 (the corporate action as a transaction, the market data operator, and its three governing principles).

\S 9 governs the one place the closed system is not closed: the boundary where external
data crosses in. Everything inside the ledger conserves quantity by construction
(Chapter~\ref{ch:objects}); value, by contrast, is supplied from outside --- prices,
dividend figures, index levels, fixings, corporate-action notices, stamped valuations.
The conservation properties hold only because the system is closed, so the boundary is
where they are won or lost. Data admitted without provenance, a corporate action that
silently re-reads old data in a new frame, a fallback taken without a record: each
corrupts the closed-system property without violating any single internal invariant. This
chapter holds the boundary watertight, and it does so with two devices --- the recorded
observation and the market data operator --- and no third.

\subsection{The boundary, and corporate actions on it}

Two things cross the boundary, and they are recorded the same way. \emph{Market and
reference data} --- a price, a dividend figure, an index level, a fixing, a stamped NAV
--- crosses as a recorded observation. A \emph{corporate action} --- a split, a dividend,
a spin-off, a merger --- crosses first as a recorded notice and is then, on its effective
date, a transaction like any other: a bundle of moves, unit-state changes, and a new
version of the instrument's terms, applied atomically through the one Transaction Executor
(Chapter~\ref{ch:machines}). A corporate action is not a privileged operation on a store;
it is admitted at the same door as a trade, and it writes ProductTerms, UnitStatus, and
balances through the same four-kind transaction (Chapter~\ref{ch:objects}).

A corporate action does one thing no trade does: it changes the meaning of all data
recorded before it. After ACME's two-for-one split, the recorded close $100.00$ still
names the same economic fact, but a post-split reader who multiplies it by the post-split
count $20{,}000$ books a phantom. To carry pre-event data into the post-event frame a
transformation must be applied, and the framework gives that transformation a name: the
\textbf{market data operator}. Three principles govern it, taken from \S 9 and used
throughout this chapter:

\begin{enumerate}[label=(\roman*),nosep,leftmargin=2.2em]
\item the operator is \emph{intrinsic to the pairing} of a kind of data with a kind of
  event, declared once when the kind of data is introduced and the corporate action is
  recorded, never improvised at the moment of reading;
\item the \emph{ledger is the authority} and the pricing data layer merely applies: the
  operator is declared data on the record, and the pricing data layer computes with it,
  never around it;
\item \emph{originals are never overwritten}. The observed data stays on the record
  exactly as observed; the adjusted value is computed at each read; derived quantities are
  recomputed from operator-adjusted inputs --- until fresh post-event data is recorded,
  after which the operator has nothing left to adjust.
\end{enumerate}

\subsection{The observation door is the one door}\label{sec:obs-door}

No data enters the ledger as a bare read, and no data enters through a door of its own.
All data crosses as a recorded \textbf{observation} carrying its provenance --- source,
observation time, basis frame. An untrusted ingestion contract \emph{proposes} a moveless
\textbf{observation-recording transaction} --- empty move list, folding the observation
into a home (Chapter~\ref{ch:homes}) --- and the one Transaction Executor admits or
refuses it like any transaction. \emph{The record is the log}: no separate observation
store, and the single-writer discipline of \S5 reaches every observation. This is the door
Chapters~\ref{ch:valuation} and~\ref{ch:virtual} name when a mark ``enters through the
observation door.'' A bare read against a live feed leaves no record and is unreproducible;
a recorded observation is a fact of the log, so every downstream number is a function of
the record alone (\S 12).

\begin{lstlisting}
data Provenance  = Provenance { source :: SourceId, observed :: Time, basis :: BasisRef }
data Observation = Observation { reading :: Reading, prov :: Provenance }
-- all data crosses as a recorded Observation; there is no bare read, and no second door.
-- an untrusted ingestion contract PROPOSES a moveless observation-recording transaction;
-- the one Transaction Executor admits or refuses it, like any transaction.
-- admission writes a moveless Transaction (empty move list) that folds the Observation
-- into a home; refusal is a returned value, never a partial write:
data Ingest      = Proposed Transaction   -- moveless: folds the Observation into a home
                 | Refused  IngestError
\end{lstlisting}

The provenance frame is what makes ``prices and quantities are never mixed across an
event'' (\S 8) enforceable rather than hoped for: the frame the data was observed in is on
the record beside the data, so a read that would combine it with a quantity in a
different frame is detectable, not silent. That detection is the invariance witness of
Section~\ref{sec:witness}.

\medskip\noindent\textbf{Episode --- the variance swap (T5): the fixing source is the
record.}
\emph{Trigger.} VS-1's scheduled fixing dates arrive; each is a watch on the recorded IDX
official close (Chapter~\ref{ch:objects}). \emph{The door.} The fixing is not a second
feed; it is a read of the same recorded official-close observation that FUT-IDX settles
against. The three settlement days of the future --- IDX closes $1{,}000.00$, $990.00$,
$1{,}005.00$ --- are the closes $C_{124}, C_{125}, C_{126}$ read at fixings $124$,
$125$, and $126$ of the swap's $252$-fixing series, each fixing being the log-return
over the prior close (fixing $125$ reads $\ln(990/1000)$, fixing $126$ reads
$\ln(1005/990)$).
\emph{Design point.} \emph{One recorded observation source feeds two consumers; because
the source is on the record, the fixing series is deterministic and replayable at the data
boundary, and firing $127$ recomputes the same statistic any party would
(Chapter~\ref{ch:valuation}).}

\medskip\noindent\textbf{Episode --- the total return swap (T6): the bridge is an
observation.}
\emph{Trigger.} TRS-1's first quarterly reset arrives; virtual ledger V-1's replayable
level is stamped at \USD{10{,}300{,}000}. \emph{The door.} That level crosses as an
observation with provenance --- source the fold of V-1's own log, observed at the reset
date, in the reference-currency frame --- and is recorded. No move crosses: a move from a
wallet of V-1 to a wallet of the true ledger has no paired leg that conserves in either
ledger, so neither Transaction Executor can admit it (Chapter~\ref{ch:virtual}). \emph{The
door checks.} Provenance present; the crossing is a recorded observation, not a move.
\emph{Design point.} \emph{The bridge between the virtual ledger and the true ledger is an
observation, never a move; the reset and financing moves it triggers are the true ledger's
own (Chapter~\ref{ch:virtual}).}

\subsection{The data-kind registry}\label{sec:registry}

Market data is not one number per name. A dividend forecast is a composite: cash entries
that carry a price dimension and scale under a split, and proportional entries that carry
no dimension and stand. An operator that treated the data as a single scalar would either
scale the proportional part, which is wrong, or leave the cash part, which is also wrong.
The operator must therefore know the \emph{dimension of each component}, and it learns it
from the \textbf{data-kind registry}.

A kind is registered once, before any data of that kind crosses, with a component-wise
declaration: each component's dimension --- price-in-basis, proportional, or quantity ---
and the quotation convention it renders. Registration is immutable: a correction is a new
kind plus a retirement of the old, never an edit, so no data ever changes dimension under
a reader's feet. Data of an unregistered kind is not silently framed: the arrival is
quarantined, its processing refused (W4, Section~\ref{sec:failures}), never dropped.

The registry has a second column under the same immutable discipline: the \textbf{event
kind} --- the kind of thing that fires a contract: a split, a fixing, a fill, an external
fail notice. The kinds of event the system can process form a finite, closed, declared set
(constitution C-4.12, ratified 2026-07-14); a kind is registered, on the record, before any
event of it crosses, and registration names its \textbf{routing} --- the contract or
contracts that must fire on it, total at registration: a kind cannot register without a
router. ``Fires the responsible contract'' (Chapter~\ref{ch:machines}) is thereby a
declared fact, not an assumption. One cause may route to many units --- one ACME dividend,
one firing per referencing unit --- separated by the cause-derived identifier. Closed and
declared, the event-kind set is the alphabet of the paths that could have occurred: what
makes the generated-event universe and the simulation sample space well defined (C-2.8).

\begin{lstlisting}
data Dim      = Price | Proportional | Quantity          -- component dimensions
data Reading  = Reading { kind :: KindId, comps :: [(Dim, Integer)] }  -- exact minor units
data IngestError = UnknownUnit | UndeclaredConvention | UnregisteredKind
data Routing  = Routing { eventKind :: EventKind, fires :: NonEmpty ContractId }
-- second column: every registered event kind names the contract(s) that must fire;
-- a kind cannot register without a router --- routing is total at registration.
-- the market data operator acts on Price components; Proportional stands; Quantity is
-- walled off to the move plane (Chapter 3) and is never carried on the value frame.
\end{lstlisting}

The wall between \texttt{Quantity} and the \emph{market data} operator is the same wall as
the move plane. A share count is not adjusted by the market data operator: the two-for-one
split moves $10{,}000$ ACME to $20{,}000$ by a paired-leg move from the issuer's virtual
wallet (Chapter~\ref{ch:objects}), because quantity is conserved and only a move may change
it. The market data operator acts on value-frame components only. This division is what makes
the split's arithmetic come out invariant, and it is stated as the witness next. (One
operator does carry a \texttt{Quantity}-dimensioned thing --- the term-adjustment operator
acting on a contract's multiplier, Section~\ref{sec:term-operator} --- and there the wall is
not breached, because a multiplier is a term, not mass.)

The registry is a named trust assumption, and it is the only one the operators rest on.

\begin{quote}
\textbf{TA-KIND.} A kind declaration faithfully renders the publisher's quotation
convention --- which components are cash and which are proportional, in which basis.
\emph{Owner:} the party that governs kind registration (data governance). \emph{Violation:}
an operator scales a component it should not, or spares one it should scale, misvaluing
across every subsequent corporate action. \emph{Detection:} the invariance witness at the
first crossing of a corporate action (Section~\ref{sec:witness}), divergence against an
independent source (failure mode W3 below), and \S 13 recomputation. \emph{Residual:} a
misdeclaration that never meets a corporate action is invisible from inside the boundary
--- a reconciliation matter at the perimeter, the exact sibling of the collateral regime
bit of Chapter~\ref{ch:collateral}.
\end{quote}

\subsection{The market data operator and adjustment-schedule totality}\label{sec:operator}

An operator is a function of the registered kind and the event kind, declared once and
recorded. It carries value-frame data from the pre-event basis to the post-event basis;
it is the identity on proportional components and undefined on quantity components, which
never reach it.

\begin{lstlisting}
-- one operator per (event kind, data kind), declared once, never improvised at read:
operator :: EventKind -> KindId -> (BasisFrame -> BasisFrame)
\end{lstlisting}

The first mechanism is that this declaration is \emph{total}.

\begin{principle}[Adjustment-schedule totality]\label{prin:sched-total}
For every registered data kind and every event kind the boundary admits, an operator is
declared --- the identity operator where a kind is genuinely unaffected, but declared, not
assumed. There is no (data kind, event kind) pair for which a read finds the adjustment
undefined. A corporate action whose pairing has no declared operator is not processed: the
arrival is recorded as a quarantined event and its adjustment refused (W4), never silently
treated as identity.
\end{principle}

Totality is what turns ``never improvised at read'' into a checkable property. An
improvised adjustment is precisely a pairing with no declared operator; forbidding the
improvisation means requiring the declaration in advance, for the whole registered domain.
The refusal of an undeclared pairing is failure mode W4 (Section~\ref{sec:failures}), and
it fails closed, as \S 2's Correctness commitment demands: the system stops rather than
guesses a frame.

Operators compose. Two corporate actions in sequence adjust data by the composition of
their operators, and the composition is taken \emph{in log order}, because the operators do
not commute: a split then a special dividend is not the same adjustment as the dividend
then the split. The log fixes the order, so the composed adjustment is well defined and
replayable.

\begin{secondtelling}
Fix a data kind. Its bases are the objects of a small category; each recorded event
contributes the arrow its operator names, from the pre-event basis to the post-event basis;
identity events contribute identity arrows. A chain of corporate actions is then a path, and
the adjustment it induces is the composite arrow --- associative because composition of
functions is, order-sensitive because the arrows do not commute. The map sending each event
sequence to its composed operator is a functor from event paths to basis maps; adjustment
totality is the statement that this functor is defined on every arrow of the source. Nothing
below depends on this telling.
\end{secondtelling}

\subsection{The operator menu: the split}\label{sec:menu}

The split shows the whole menu at once. On 2026-06-01 ACME's two-for-one split is
effective; the recorded split notice had declared the effective-date watch, the Event
Monitor emits the event, and one transaction carries the move and the new terms version
(Chapters~\ref{ch:objects} and~\ref{ch:homes}). What each recorded data kind reads as, in
the new frame, is the operator's work:

\begin{center}
\begin{tabular}{@{}lll@{}}
\toprule
Data kind & Operator for the 2-for-1 split & Read after the event \\
\midrule
Official close (spot) & scale by $\tfrac12$ & $100.00 \rightarrow 50.00$ \\
Cash-per-share dividend figure & scale by $\tfrac12$ & $1.50 \rightarrow 0.75$ \\
Proportional (percentage) figure & identity --- stands & $1.00\% \rightarrow 1.00\%$ \\
Index divisor (a constituent splits) & adjusted to hold the level & level continuous \\
\bottomrule
\end{tabular}
\end{center}

\noindent The share count is absent from the menu on purpose: $10{,}000 \rightarrow
20{,}000$ is a move, not an operator (Section~\ref{sec:registry}). The spot and the
cash-per-share figure are price-dimensioned and halve; the percentage figure is
proportional and stands; the index divisor is adjusted so the published level does not jump
when a constituent splits. Every original stays on the record; each ``read after'' is
computed at the read, until a fresh post-split close prints and the operator retires for
that data.

\medskip\noindent\textbf{Episode --- the dividend figure under the split (T3$\times$T4).}
\emph{Trigger.} On 2026-05-04 ACME's cash dividend of \USD{1.50} per share is announced;
the announcement is captured at the boundary as a recorded observation, and its declared
per-share figure $1.50$ enters the record (the three firings that pay it are
Chapter~\ref{ch:contracts}'s). \emph{The operator.} The split is effective 2026-06-01,
after the payment date but within the window in which the declared forward figure is still
read; a read of the $1.50$ figure in the post-split frame applies the split's
cash-per-share operator and returns $0.75$. \emph{The door checks.} The data kind is
registered as a cash-per-share (price) component, so the operator scales it; a proportional
figure beside it would have stood. \emph{Design point.} \emph{One operator, one invariant,
visible: the same halving that takes the spot from $100.00$ to $50.00$ takes the recorded
per-share dividend figure from $1.50$ to $0.75$, so the dividend history reads in
current-share terms --- the scaling of a real figure is shown, not asserted.}

\subsection{The term-adjustment operator: a derivative's own terms}\label{sec:term-operator}

The menu above adjusts \emph{value-frame} data --- a spot, a dividend figure, an index
level. A unit whose terms \emph{reference} the split stock has a second frame the same event
must adjust: its own \textbf{declared terms}. OT-1's barrier of $80.00$ names OMEGA's
official close; when OMEGA splits two-for-one every OMEGA price halves, and a barrier left at
$80.00$ knocks on the first post-split close near $47.50$ --- a stock at half its former level
meeting a barrier struck against the whole. The barrier is a declared term, not data, so
no operator of the menu touches it. The framework gives the missing transformation the exact
sibling of the market data operator: the \textbf{term-adjustment operator}. Where the market
data operator carries \emph{data} from a pre-event basis to a post-event basis, the
term-adjustment operator carries a \emph{declared term} across the same event; it is the
action of the corporate-action out-edge every referencing unit declares
(Chapter~\ref{ch:objects}, Section~\ref{sec:termsheet-graphs}), which appends a new
ProductTerms version with each referencing term carried through its operator.

\begin{lstlisting}
-- the sibling of  operator  (Section 8.4), acting on the declared-term frame, not value:
termOp :: EventKind -> TermKind -> (Term -> Term)   -- one per (event kind, term kind)
\end{lstlisting}

It is total in exactly the sense of Principle~\ref{prin:sched-total}: for every
(declared-term kind, event kind) pair an operator is declared --- the identity where a term
is genuinely unaffected, but declared, not assumed --- and a pairing with no declared operator
refuses the event at the door. Totality here is what forbids a silent stale barrier: an
undeclared (term kind, event kind) pairing cannot cross, so a split can never leave a
referencing term unadjusted by default.

The operator reuses the dimension machinery of the registry (Section~\ref{sec:registry}): a
term carries a dimension, and the dimension fixes how the event moves it.

\begin{center}
\begin{tabular}{@{}llll@{}}
\toprule
Declared term & Dimension & Operator for the 2-for-1 split & Read after \\
\midrule
Barrier (OT-1 on OMEGA)          & \texttt{Price}        & scale by $\tfrac12$             & $80.00 \rightarrow 40.00$ \\
Strike (an OMEGA option)         & \texttt{Price}        & scale by $\tfrac12$             & $K \rightarrow K/2$ \\
Multiplier (a single-name        & \texttt{Quantity}     & scale by $2$ --- \emph{doubles} & $100 \rightarrow 200$ \\
\quad contract on OMEGA)         &                       &                                & \\
Participation rate               & \texttt{Proportional} & identity --- stands             & $150\% \rightarrow 150\%$ \\
Fixed cash payout (OT-1)         & declared absolute     & identity --- stands             & \USD{1{,}000{,}000} unchanged \\
\bottomrule
\end{tabular}
\end{center}

Note the direction. A \texttt{Price}-dimensioned term --- a barrier, a strike --- halves,
exactly as the spot halves, because it is quoted in the units the price is quoted in. A
\texttt{Quantity}-dimensioned term --- the multiplier of a single-name contract on the split
stock --- \emph{doubles}, because the split doubles the number of underlying shares one
contract commands and the payoff $(\text{level} - K)\times\text{multiplier}$ must stand:
$(S - K)\times 100 = (\tfrac{S}{2} - \tfrac{K}{2})\times 200$. A \texttt{Proportional} term
--- a participation rate --- stands. And a fixed cash payout is neither \texttt{Price} nor
\texttt{Proportional} but a \emph{declared absolute}: it is registered as such, and its split
operator is the identity --- \USD{1{,}000{,}000} is \USD{1{,}000{,}000} whether OMEGA is whole
or halved.

The multiplier is the one place a \texttt{Quantity}-dimensioned thing is carried by an
operator. On the balance plane the wall of Section~\ref{sec:registry} is absolute: only a
move changes a quantity, and no market data operator touches a share count. The term-adjustment
operator carries a \texttt{Quantity}-dimensioned multiplier and doubles it, and this does not
breach that wall, because \emph{a multiplier is a term, not mass}. It names how many
underlying units one contract controls; adjusting it re-describes the same contract in the
post-split frame and moves no holding of anything. The share count the split really moves ---
$10{,}000$ OMEGA to $20{,}000$ --- still travels by a paired-leg move and by nothing else
(Section~\ref{sec:registry}). Mass is carried by moves; a term is carried by its operator; the
value frame and the balance plane never touch.

\paragraph{Declared identity is still a declaration: VS-1 and FUT-IDX on the index.} Not every
referencing unit adjusts under every event, but totality requires the non-adjustment be
\emph{declared}. VS-1's $252$ fixings read the $253$ IDX official closes $C_0,\dots,C_{252}$, and FUT-IDX settles the IDX
level at multiplier~$10$; both reference IDX, so both declare a corporate-action out-edge. When
a \emph{constituent} of IDX splits, the index divisor is adjusted to hold the published level
continuous (the menu's divisor row above), so the level VS-1 and FUT-IDX read does not jump.
Their term-adjustment operator for the (IDX-constituent split, fixing) and (IDX-constituent
split, multiplier) pairings is therefore the \emph{identity} --- but declared, per
Principle~\ref{prin:sched-total}, not assumed: FUT-IDX's multiplier $10$ stands because the
level it multiplies is held continuous by the divisor, and VS-1's fixing reads the same
continuous level. The contrast with the single-name multiplier is the whole point: a contract
on OMEGA \emph{itself} doubles its multiplier when OMEGA splits, because OMEGA's price halves;
a contract on the IDX \emph{level} leaves its multiplier alone when a constituent splits,
because the divisor already held the level. Both are declared operators; neither is silent.

\subsection{State-aware reads: consistency of reference and cum/ex}\label{sec:witness}

Every operator read is state-aware: the frame data reads in is determined by the
corporate-action state of its underlying at the read, and that state is on the record. Two
readings of one recorded observation in two frames are two different values, and combining a
quantity with a price drawn from a different frame produces a phantom. The check that
refuses phantoms is the \textbf{invariance witness}: the operator on the value frame and
the move on the quantity are bound so their product --- total value --- is invariant across
the event.

For the split the witness is arithmetic. Before: $10{,}000 \times 100.00 = \USD{1{,}000{,}000}$.
After, read in the post-split frame: $20{,}000 \times 50.00 = \USD{1{,}000{,}000}$. The
phantom the witness refuses is $20{,}000 \times 100.00 = \USD{2{,}000{,}000}$, the new count
against the stale frame. This is the consistency-of-reference invariant; it is stated
formally and catalogued in Chapter~\ref{ch:invariants}, and used here.

For a derivative the witness has no move to bind, and it is stated over the declared terms
instead. OT-1 references OMEGA but carries no OMEGA quantity: its holder holds one option
before OMEGA's split and one after, so no move fires. The witness is then that the
\emph{payoff} on the post-event term frame equals the payoff on the pre-event term frame.
Under OMEGA's two-for-one split the term-adjustment operator (Section~\ref{sec:term-operator})
carries OT-1's barrier $80.00 \rightarrow 40.00$ and leaves its \USD{1{,}000{,}000} payout
unchanged --- a declared absolute stands. The witness is the knock predicate read in the new
frame: a pre-split close at or below $80.00$ pays \USD{1{,}000{,}000}, and a post-split close
at or below $40.00$ pays the same \USD{1{,}000{,}000} --- one economic event, two frames, one
payout. The phantom the witness refuses is the dropped-edge failure of Chapter~\ref{ch:objects}
(Section~\ref{sec:termsheet-graphs}): the first post-split close prints $\approx 47.50$;
against the stale barrier $80.00$ it reads $47.50 \le 80.00$ and knocks, paying
\USD{1{,}000{,}000} on a stock that never fell; against the adjusted barrier $40.00$ it reads
$47.50 > 40.00$ and no knock fires. The stale-frame read is the phantom; the term-adjustment
operator, appended as the corporate-action edge's new terms version, is what refuses it.
Quantity-times-price for a holding and payoff-on-terms for a derivative are the two readings of
one witness.

The same witness governs the cum/ex boundary of a distribution, where no quantity changes
at all. Before the ex date a price is \emph{cum} --- it includes the coming distribution;
on and after the ex date it is \emph{ex} --- it excludes it. The cum price and the ex price
are one data kind read in two states separated by the ex-date event, whose declared
operator subtracts the distribution's value. A cum price consumed as though it were ex, or
an ex price consumed as though it were cum, is a phantom of exactly the split's kind ---
the right number in the wrong frame --- and the same invariance witness refuses it, because
the state that fixes the frame is the recorded ex-date event, not the wall clock. A read
that ignores the ex-date state prices the instrument on the wrong side of its own corporate
action, and the witness catches it before the number reaches a valuation.

One thing the ex-date does not carry is a promise about \emph{when} an ex trade settles
relative to the record date. That the ex-date is set so an ex trade settles only after the
record date --- keeping the registered holder and the entitled party the same party --- is
an external market convention, fixed by the corporate-action calendar and the settlement
cycle, and it is nowhere on the ledger's record. The ledger records the ex-date as data
with provenance and routes entitlement against it (Chapter~\ref{ch:settlement}); it does not
enforce the convention that keeps registration and entitlement aligned. That convention
breaks lawfully under a due-bills special dividend, whose ex-date falls after the payment
date and lets an ex trade settle before the record date. The coherence is named as the
trust assumption TA-EXDATE in Chapter~\ref{ch:requirements}, reconciled against at this
boundary but not enforced here; the settlement interface handles the case it breaks with the
mirror market-claim, so correctness does not rest on the convention holding.

\subsection{Recompose and the aggregate weld}\label{sec:recompose}

The split adjusts one unit's data in place. Two further mechanisms handle corporate actions
that change which units exist: one splits a unit into several, the other merges several into
one. Both preserve total value, and both are checked by the same witness.

\paragraph{Recompose (one unit becomes many): a spin-off.}
When an issuer distributes a new entity to its holders, one position becomes a basket, and
the pre-event data of the parent must be carried onto the resulting units. \emph{Recompose}
is the operator that does so, per a declared allocation basis recorded with the event.
Illustratively: HELIOS trades at $120.00$ and spins off LUMEN, one LUMEN share per HELIOS
share, with a declared allocation of $75\%$ of pre-event value to the HELIOS stub and
$25\%$ to LUMEN. A holder of $1{,}000$ HELIOS holds, after the event, $1{,}000$ HELIOS and
$1{,}000$ LUMEN --- the new LUMEN mass arriving by a paired-leg move from the issuer, not by
an operator. Recompose carries the recorded pre-event close onto the two frames: HELIOS stub
$90.00$, LUMEN $30.00$. The witness holds: $1{,}000 \times 120.00 = 1{,}000 \times 90.00 +
1{,}000 \times 30.00 = \USD{120{,}000}$. The original $120.00$ is preserved; the stub and
spun-off frames are computed at read until fresh post-event closes print for each.

\paragraph{The aggregate weld (many units become one): an elective merger.}
When one unit is absorbed into another, the resulting position welds into the acquirer's
aggregate, and the two data histories must weld into one continuous series so that valuation
and profit and loss run unbroken across the event. The \emph{aggregate weld} is the operator
that carries the pre-event data of the absorbed unit onto the acquirer's frame, per the
elected terms recorded at the boundary. An elective merger is captured as an observation ---
the election result is data with provenance, like any other. Illustratively: TARGET is
acquired by ACQUIRER, each TARGET share electing either \USD{50.00} cash or $0.4$ ACQUIRER
shares; ACQUIRER trades at $125.00$. A holder of $1{,}000$ TARGET who elects stock welds to
$400$ ACQUIRER, and the pre-merger TARGET series is welded onto the ACQUIRER frame at the
elected ratio $0.4$; the witness holds, $1{,}000 \times 50.00 = 400 \times 125.00 =
\USD{50{,}000}$. A holder who elects cash has the position retired at settlement and no weld
occurs --- the elected terms, recorded data, decide which mechanism runs.

Two further composite distributions --- a dividend forecast and a special dividend --- are
governed by the same registry and the same operators: the forecast is a composite of cash
and proportional components that the split scales component-wise, and the special dividend
carries a cum/ex operator like the ordinary one. Their fully worked arithmetic is deferred
to the Worked-Examples companion (register lines~E71 and~E72); the mechanism here is
complete without them.

\subsection{Failure modes at the boundary}\label{sec:failures}

The boundary has four failure modes W1--W4, a closed set distinct from the three collateral
inflow \emph{regimes} of Chapter~\ref{ch:collateral}: two closed enumerations, two names, never
one word for both. Each mode has a determined, replayable disposition --- a
function of log state alone, so replay reproduces it bit-identically. None improvises; each
fails closed or records a flag, never a silent pick.

\begin{center}
\begin{tabular}{@{}lll@{}}
\toprule
 & Boundary condition & Determined, replayable disposition \\
\midrule
W1 & missing data & the operator applies to the last recorded original, carrying a \\
   &               & staleness flag; with no original, the price is undefined, not \\
   &               & zero, and valuation over the unit defers (Chapter~\ref{ch:valuation}) \\
W2 & late data & recorded at its observation-time coordinate; values derived over \\
   &            & $[\text{observed},\text{arrival})$ are recomputed --- ``as known at $t$'' and \\
   &            & ``with corrections through $t'$'' are both queryable \\
W3 & contradictory sources & sources diverging beyond the declared threshold yield a \\
   &                       & flagged ``aggregation failed'' observation, never a silent \\
   &                       & pick; sources and threshold are recorded \\
W4 & unroutable or & the arrival is recorded as a \textbf{quarantined event}, a \\
   & unmappable event & moveless provenance-carrying fact, and its \emph{processing} \\
   &                  & is refused --- the refusal on the record; fail-closed, nothing lost \\
\bottomrule
\end{tabular}
\end{center}

\noindent W1 is the working state of \S 9's ``until fresh post-event data is recorded'':
between the event and the first fresh observation, every read is the original under its operator,
flagged stale, and the flag is a fact of the record, not a property of a cache. W2 makes the
boundary bitemporal, and it is the machinery Chapter~\ref{ch:picture}
(Section~\ref{sec:bitemporal}) and Theorem~\ref{thm:timetravel} rest on. Every observation
carries two coordinates: its \emph{valid-time} coordinate --- its observation time, or the
fixing index~$k$ it belongs to --- and its \emph{transaction-time} coordinate --- the log
position at which it was recorded. Data arriving late does not rewrite history: it is
appended at its own transaction-time position while keeping its earlier valid-time
coordinate, so the as-known-at view --- the projection over any prefix that predates the
arrival --- is preserved, and the as-of view --- the projection over the prefix that now
carries the correction --- recomputes the derived folds over the affected valid-time
interval, bit-exactly (\S 2). Both views survive; neither overwrites the other. An amended
or withdrawn corporate-action notice is this same case --- superseding data recorded at
its own transaction-time position, never an edit of the original notice, so the
pre-amendment frame and the corrected one both survive. C-12.1 (v1.2) requires both honest
answers --- as known at $t$, and with corrections through $t'$ --- and delegates their
indexing to this specification; W2 is where the two-axis indexing discharges that duty
(entry PARK-2 closed, Chapter~\ref{ch:scope}). W3 is the
multi-source case: a single source is a single point of
failure, so where data materially affects valuation it is drawn from more than one
attested source, aggregated by a declared rule, and a divergence beyond the declared
threshold produces a flagged output rather than a silently chosen one. The divergence
threshold is a governance parameter, recorded and calibrated, not a silent constant: under
correlated stress --- every venue diverging at once --- it defers valuation book-wide rather
than choose a number, an intended fail-closed whose calibration is a governance matter, not a
hidden behaviour. W4 is
totality enforced without loss. An event of an undeclared kind, an unresolvable reference,
or a pairing with no declared operator (Principle~\ref{prin:sched-total}) is neither
silently dropped nor silently guessed: it crosses the one door as a quarantined event ---
a moveless recorded fact carrying its provenance --- and its \emph{processing} is refused,
the refusal itself a fact of the record. The system still stops rather than improvises
(\S2); it no longer forgets what it stopped on. The quarantined event mints a visible,
deadline-bearing open item and stands until its kind is registered or a person disposes of
it (C-12.4 phase-0, constitution C-4.12).

When the kind is later registered, or a person disposes, the arrival is processed then: the
processing transaction's cause is the recorded quarantined event --- idempotent under the
cause-derived identifier, carrying the quarantined event's own valid-time coordinate (W2
above). The registry is versioned declared data read in force at each log position, so
replay across its evolution is deterministic: every reader computes one answer.

\subsection{The boundary-convention slot}\label{sec:slot}

One class of boundary fact is neither data nor an operator: a \emph{convention} that
some downstream projection needs and that must be stated once, as data, so that two readers
of one log compute one answer. \S 9 places such conventions on the record beside the
operators; this section declares the \textbf{slot} they occupy, and
Chapter~\ref{ch:reporting} fills and consumes one.

A convention in the slot is a total function of log state. The convention
Chapter~\ref{ch:reporting} consumes --- the attribution of commingled received value across
a reinvestment pool --- has three components, and the slot's shape is exactly these three:

\begin{lstlisting}
data Convention = Convention
  { pool     :: Predicate (Wallet, Move)   -- the pool over which value is attributed
  , basis    :: PeriodBasis                -- when the pro-rata ratio is struck
  , rounding :: RoundingRule }             -- remainders booked explicitly (Chapter 3)
attribute :: Convention -> LogPrefix -> Partition Amount   -- consumed in Chapter 12
\end{lstlisting}

The slot declaration is itself a recorded event, so it versions with the log prefix: a
report computed as of a past date uses the convention then in force, and a changed
convention is a new recorded version, never an edit. Totality is the requirement that makes
the slot safe to consume: for every log prefix the convention assigns every attributable
quantity, the attributed parts and the explicit remainder partition the whole, and the
rounding rule leaves no quantity silently dropped (\S 4's remainder discipline). A
convention missing any component is refused, because two readers would then compute
different figures from one log --- internal reconciliation failure reintroduced through a
report. Chapter~\ref{ch:reporting} states the totality property and proves the figures it
builds are projections of coordinates, units, and this declared convention alone; this
chapter's duty is discharged by declaring the slot and its totality obligation.

\subsection{What an auditor checks}

A candidate implementation satisfies this chapter when seven things hold, each mechanically
checkable over generated data and histories: all data on the record carries provenance,
and no valuation reads data that does not (Section~\ref{sec:obs-door}); every data kind
is registered with component dimensions before it crosses, and an unregistered kind is
refused (Section~\ref{sec:registry}); an operator is declared for every (registered kind,
event kind) pair, and an undeclared pairing is refused, not defaulted to identity
(Principle~\ref{prin:sched-total}); every original is preserved and every adjusted value is
recomputed at read from the original under the declared operators
(Sections~\ref{sec:menu}--\ref{sec:recompose}); the invariance witness holds across every
generated corporate action, and a stale-frame or wrong-side-of-ex read is refused
(Section~\ref{sec:witness}); each of W1--W4 produces its determined disposition on
replay (Section~\ref{sec:failures}); and every unroutable or unmappable arrival is recorded
as a quarantined event, never dropped, and re-processes deterministically once its kind is
registered (Sections~\ref{sec:registry},~\ref{sec:failures}). These restate as minima in
Chapter~\ref{ch:requirements}. The boundary is then closed: what crosses it is recorded,
framed, and replayable, and the conservation the rest of the specification proves rests on
solid ground.
